\chapter[Conclusão]{Conclusão}

Este trabalho investigou como uma arquitetura reprodutível de Engenharia de Dados pode integrar bases educacionais heterogêneas do INEP e disponibilizá-las de forma consistente para análise no Power BI. A questão de pesquisa foi respondida por meio da construção de um pipeline em Python organizado nas camadas Raw, Bronze, Silver e Gold, acompanhado de manifesto das fontes, hashes SHA-256, metadados de linhagem, validadores independentes e uma camada analítica consumida exclusivamente pelo Power BI.

O pipeline integrou IDEB, SAEB, Rendimento Escolar e TDI no recorte histórico de 2007 a 2023 e incorporou a PND 2025 como análise complementar. A solução final materializou cinco tabelas fato e cinco dimensões. A validação global confirmou a integridade estrutural do modelo e as correspondências definidas entre as camadas, retornando \texttt{MODELO DIMENSIONAL GOLD: OK}.

A construção mostrou que a integração de fontes oficiais exige mais do que reunir arquivos em uma ferramenta de visualização. Foi necessário tratar diferenças de formato, layout, periodicidade, granularidade e semântica de rede, além de documentar decisões específicas por indicador. No SAEB, por exemplo, a preservação dos agregados oficiais de UF evitou substituir o valor publicado por uma recomposição escolar que não reproduzia o resultado oficial. Na PND, a auditoria da população e o uso dos pontos de corte aplicados à \texttt{NT\_OBJ} substituíram critérios analíticos anteriores e produziram uma classificação coerente com a documentação adotada no projeto.

O Power BI atuou como camada de consumo e análise. As 27 medidas DAX e as páginas do dashboard permitiram observar panorama histórico, desempenho por Unidade Federativa, associação descritiva entre aprendizagem e fluxo, variações territoriais, distorção idade-série e resultados da PND 2025. As medidas temporais foram implementadas de forma a usar o primeiro e o último ano com dado dentro do intervalo selecionado, respeitando a periodicidade anual de Rendimento e TDI e a periodicidade bienal de IDEB e SAEB.

Os resultados descritivos evidenciam trajetórias distintas entre indicadores e UFs. No panorama geral do modelo, observam-se aumento do IDEB, das proficiências e da aprovação ao longo de grande parte da série, acompanhado de redução da TDI. Na página específica de TDI, o intervalo 2007--2023 apresenta variação média de -18,3 pontos percentuais no contexto exibido, enquanto a PND reúne 759.140 participantes e 493.208 classificados como proficientes, correspondentes a 64,97\% da população analítica. Esses resultados são descritivos e não demonstram relações causais.

A principal contribuição técnica do trabalho está na reprodutibilidade do percurso entre fonte e consumo analítico. O inventário de 54 arquivos Raw, a identificação por hashes, a documentação da origem, os validadores de cada camada, o orquestrador e os testes do repositório tornam explícitas decisões que, em um fluxo restrito à ferramenta de BI, poderiam permanecer ocultas. Os notebooks demonstrativos complementam essa documentação, sem substituir os scripts oficiais do pipeline.

As limitações permanecem relevantes. A série histórica é harmonizada principalmente em nível de UF, rede pública e etapas agregadas, o que impede inferências sobre escolas ou municípios. As fontes possuem metodologias e layouts que mudam ao longo do tempo; SAEB apresenta particularidade de rede em 2007 e 2009; IDEB e SAEB são bienais; a PND possui apenas 2025; e as médias do dashboard não equivalem automaticamente a médias nacionais ponderadas. Além disso, as visualizações não permitem inferência causal.

Dessa forma, conclui-se que a arquitetura desenvolvida atingiu o objetivo proposto: bases públicas heterogêneas foram integradas em um pipeline documentado, validado e rastreável, chegando a um modelo dimensional consistente para exploração no Power BI. A contribuição do trabalho está tanto no produto analítico final quanto na explicitação do processo de Engenharia de Dados necessário para produzi-lo.

\section{Principais Contribuições}

As principais contribuições podem ser sintetizadas em cinco eixos. Primeiro, a integração de fontes educacionais heterogêneas do INEP em um fluxo único. Segundo, a arquitetura em camadas, que separa preservação, ingestão, harmonização e consumo analítico. Terceiro, a rastreabilidade e a reprodutibilidade, apoiadas por manifesto, hashes, linhagem, validadores, testes e orquestração. Quarto, a modelagem dimensional com cinco fatos e cinco dimensões. Quinto, a disponibilização dos dados em um modelo semântico com 27 medidas DAX e dashboards temáticos.

No componente territorial, a padronização histórica por Unidade Federativa viabilizou a integração entre fontes com diferentes estruturas. Para a PND, o modelo preservou também o código do município de aplicação e a relação ativa com \texttt{DIM\_MUNICIPIO}, sem criar uma relação direta entre essa dimensão e \texttt{DIM\_UF}, evitando ambiguidade de filtro.

A incorporação da PND 2025 demonstrou a possibilidade de integrar uma fonte de granularidade individual ao mesmo ambiente analítico sem tratá-la como continuação artificial da série histórica da Educação Básica. A separação entre a análise histórica e a PND preservou a interpretação de cada fonte.

\section{Trabalhos Futuros}

Como trabalhos futuros, sugere-se ampliar o modelo com outras bases do INEP, como Censo Escolar, infraestrutura, matrículas, turmas e formação docente. A inclusão dessas fontes pode permitir análises adicionais, desde que cada nova base passe pelo mesmo processo de auditoria, documentação de fonte, harmonização e validação.

Outra possibilidade é aprofundar a granularidade territorial ou institucional. O presente trabalho prioriza UFs na análise histórica; pesquisas futuras podem desenvolver recortes municipais ou escolares quando houver compatibilidade de fonte, grão e metodologia.

Também podem ser incorporados indicadores socioeconômicos de instituições como IBGE ou IPEA. Essa expansão permitiria explorar associações mais amplas entre contexto social e indicadores educacionais, sem confundir análise descritiva com inferência causal.

Em estudos futuros, técnicas estatísticas ou de ciência de dados podem ser utilizadas para correlação, regressão, agrupamento ou séries temporais. Caso o objetivo seja investigar causalidade, será necessário um desenho específico de pesquisa, com hipóteses, controles e métodos adequados.

No caso da PND, novas edições poderão permitir a construção de uma série própria, separada dos indicadores históricos deste trabalho. Também é possível evoluir a publicação do dashboard para ambiente institucional ou público, com atenção a usabilidade, atualização automatizada e governança dos dados.

\section{Contribuições em Produção Bibliográfica}

Até o momento da finalização deste trabalho, não houve publicação de artigo científico resultante diretamente da pesquisa. O projeto, entretanto, oferece material para futuras produções sobre pipelines reprodutíveis de dados educacionais, qualidade e linhagem de dados públicos, modelagem dimensional aplicada ao INEP, integração de indicadores da Educação Básica e visualização analítica em Power BI.

Possíveis trabalhos derivados incluem estudos sobre reprodutibilidade de pipelines educacionais, comparação de metodologias de publicação entre edições do SAEB, integração de indicadores históricos com novas fontes e uso de visualizações para comunicação de resultados educacionais.
